% Generated draft: requires mathematical and source review.
\documentclass[11pt]{article}
\usepackage[margin=25mm]{geometry}
\usepackage{fontspec}
\setmainfont{Latin Modern Roman}
\usepackage{amsmath,amssymb,mathtools,unicode-math}
\setmathfont{Latin Modern Math}
\usepackage{enumitem,needspace,xurl,hyperref}
\setlength{\parindent}{0pt}
\setlength{\parskip}{5pt}
\setlength{\emergencystretch}{3em}
\newcommand{\N}{\mathbb{N}}
\newcommand{\Z}{\mathbb{Z}}
\newcommand{\Q}{\mathbb{Q}}
\newcommand{\R}{\mathbb{R}}
\newcommand{\C}{\mathbb{C}}
\newcommand{\F}{\mathbb{F}}
\newcommand{\A}{\mathbb{A}}
\newcommand{\PP}{\mathbb{P}}
\newcommand{\E}{\mathbb{E}}
\newcommand{\eps}{\varepsilon}
\newcommand{\dd}{\,\mathrm{d}}
\begin{document}
\section*{UnsolvedMath ID: 30004726}
\textbf{Source problem code:} catalogue entry
\subsection{Definitions and mathematical statement}
For a smooth three-dimensional manifold \(M\), let \([M]_{\mathrm{diff}}\) denote its diffeomorphism class. A smooth Riemannian metric \(g\) has pointwise positive scalar curvature if its scalar curvature function satisfies
\[
\operatorname{Scal}_g(x)>0\qquad\text{for every }x\in M.
\]
For a manifold without boundary, completeness means completeness of the Riemannian distance \(d_g\), equivalently geodesic completeness. The metric satisfying the scalar curvature inequality must itself be complete. No uniform positive lower bound on scalar curvature, or prescribed rate of decay at infinity, is required in the supplied question.

The classification problem is to give necessary and sufficient conditions on the diffeomorphism class of \(M\) for the existence of such a metric; equivalently, to describe
\[
\left\{[M]_{\mathrm{diff}}:\ \dim M=3,\ \exists\text{ a smooth Riemannian metric }g\text{ on }M,\ (M,d_g)\text{ is complete},\ \forall x\in M,\ \operatorname{Scal}_g(x)>0\right\}.
\]
This asks for a classification of the underlying manifolds, extending the classification of the closed orientable case referred to in the source. Here closed means compact and without boundary. The supplied record does not state that closed classification or specify precisely how the extension should be presented.

The admissible category of manifolds in the displayed classification remains a source qualification: the record does not explicitly specify orientability, connectedness, or whether boundary is excluded. The completeness convention above gives the usual interpretation for manifolds without boundary; if manifolds with boundary are intended, the meaning of completeness and any boundary conditions require clarification. These scope questions are unresolved rather than additional hypotheses imposed by this draft.
\subsection{Short English statement}
Classify, up to diffeomorphism, the smooth three-manifolds that admit a complete Riemannian metric whose scalar curvature is positive at every point, extending the closed orientable classification mentioned in the source. Positivity need not have a uniform positive lower bound. The intended orientability, connectedness, and boundary conventions are not specified.
\subsection{Sources}
\begin{itemize}
\item Supplied local UnsolvedMath registry record, permanent ID \(30004726\), source problem code \texttt{catalogue entry}, in the Oberwolfach Reports Open Problems collection. The record labels the statement exact and its status partially solved; these are supplied curation labels, not independent verification. No individual proposer is identified.
\item Supplied report citation: \emph{Analysis, Geometry and Topology of Positive Scalar Curvature Metrics} (2022), Oberwolfach Reports, pp.~1579--1649, DOI \texttt{10.4171/owr/2021/30}. \url{https://doi.org/10.4171/owr/2021/30}. The record gives proposed year 2021; its report title uses 2022, while the DOI contains 2021. These details are retained as supplied. No problem-specific page is given.
\item Supplied literature reference: Wang, ``topology of three-manifolds with uniform positivity.'' \url{https://arxiv.org/abs/2212.14383}.
\item Supplied literature reference: ``Complete three-manifolds with quadratic decay.'' \url{https://link.springer.com/article/10.1007/s00208-025-03192-9}.
\item Literature assessment dated exactly \texttt{2026-08-22}: the supplied curation states that uniformly positive scalar curvature is topologically classified for open three-manifolds, that several decay-controlled cases are known, and that the unrestricted complete pointwise-positive-scalar-curvature classification remains open. This is supplied dated curation derived from the OWR record, not independently verified. Its time-sensitive status and cited claims require independent checking; the precise hypotheses of those results are not supplied.
\end{itemize}
\end{document}
